\chapter{Routing with React Router}

\section{Core Building Blocks}

\begin{lstlisting}[language=JavaScript]
// App.jsx
import { BrowserRouter, Routes, Route } from "react-router-dom";
import MainLayout from "./layouts/MainLayout";
import ProtectedRoute from "./components/ProtectedRoute";
import Login from "./pages/Login";
import Register from "./pages/Register";
import Dashboard from "./pages/Dashboard";
import TaskDetails from "./pages/TaskDetails";
import Profile from "./pages/Profile";

export default function App() {
  return (
    <BrowserRouter>
      <Routes>
        <Route path="/login" element={<Login />} />
        <Route path="/register" element={<Register />} />

        <Route
          element={
            <ProtectedRoute>
              <MainLayout />
            </ProtectedRoute>
          }
        >
          <Route path="/dashboard" element={<Dashboard />} />
          <Route path="/tasks/:id" element={<TaskDetails />} />
          <Route path="/profile" element={<Profile />} />
        </Route>
      </Routes>
    </BrowserRouter>
  );
}
\end{lstlisting}

\begin{itemize}
  \item \texttt{BrowserRouter} -- enables client-side routing via History API.
  \item \texttt{Routes}/\texttt{Route} -- path-to-component mapping.
  \item \texttt{Link}/\texttt{NavLink} -- render \texttt{<a>} without reload; \texttt{NavLink} adds an active class.
  \item \texttt{useNavigate} -- programmatic navigation.
  \item \texttt{useParams} -- reads dynamic URL segments like \texttt{:id}.
  \item A path-less nested \texttt{Route} under \texttt{<Outlet/>} creates a \textbf{layout route}.
\end{itemize}

\begin{lstlisting}[language=JavaScript]
// layouts/MainLayout.jsx
import { Outlet } from "react-router-dom";
import Navbar from "../components/Navbar";
import Sidebar from "../components/Sidebar";

export default function MainLayout() {
  return (
    <div className="app-shell">
      <Navbar />
      <div className="app-body">
        <Sidebar />
        <main><Outlet /></main>
      </div>
    </div>
  );
}
\end{lstlisting}

\begin{lstlisting}[language=JavaScript]
// pages/TaskDetails.jsx
import { useParams, useNavigate } from "react-router-dom";

export default function TaskDetails() {
  const { id } = useParams();      // id from /tasks/:id
  const navigate = useNavigate();

  const handleBack = () => navigate("/dashboard");
  // ...fetch task by id, render details
}
\end{lstlisting}

\section{Protected Routes}

Protected routes stop unauthenticated users from rendering a page that needs
a logged-in user before the page component mounts.

\begin{diagrambox}[Protected Route Flow]
\begin{verbatim}
 /dashboard requested -> <ProtectedRoute> renders
   -> useAuth(): logged in?
        NO  -> <Navigate to="/login"/>
        YES -> render children (<MainLayout/> -> <Outlet/>)
\end{verbatim}
\end{diagrambox}

\begin{lstlisting}[language=JavaScript]
// components/ProtectedRoute.jsx
import { Navigate } from "react-router-dom";
import { useAuth } from "../hooks/useAuth";

export default function ProtectedRoute({ children }) {
  const { user, loading } = useAuth();

  if (loading) return <p>Checking session...</p>;
  if (!user) return <Navigate to="/login" replace />;

  return children;
}
\end{lstlisting}

\begin{notebox}[NOTE]
\texttt{replace} keeps the login page from being pushed onto history again,
so back button doesn't bounce between \texttt{/login} and the redirecting page.
\end{notebox}

\begin{warnbox}[WARNING]
A client-side \texttt{ProtectedRoute} only hides UI, not a security boundary --
Express must independently verify the JWT via its \texttt{protect} middleware
(Chapter 2) on every protected endpoint.
\end{warnbox}
